\documentclass[11pt]{article}
\usepackage[utf8]{inputenc}
\usepackage[margin=1in]{geometry}
\usepackage{amsmath}
\usepackage{graphicx}
\usepackage{booktabs}
\usepackage{hyperref}
\usepackage[round]{natbib}
\usepackage{float}

\title{Mapping Responses to Focal Injections of Bicuculline Along the Rostro-Caudal Axis of the Lateral Parafacial Region Identifies the Core Site for Active Expiration Generation}
\author{Author A$^{1}$, Author B$^{1}$, Author C$^{2}$, Author D$^{1,2}$\\[6pt]
$^{1}$Department of Cell Biology and Physiology, University of Alberta\\
$^{2}$Women and Children's Health Research Institute}
\date{}

\begin{document}
\maketitle

\begin{abstract}
The expiratory oscillator in the lateral parafacial region (pFL) of the ventral medulla lacks a definitive anatomical marker and is silent at rest, making it difficult to localize precisely. Prior pharmacological and chemogenetic studies have targeted locations spanning --0.2 to +0.5~mm from the caudal tip of the facial nucleus (VIIc), but none has systematically mapped respiratory responses across the full rostro-caudal extent. Here we injected bicuculline methochloride (200~$\mu$M, 200~nL bilateral) at five locations from --0.2 to +0.8~mm from VIIc in urethane-anesthetized Sprague-Dawley rats (n=35) and quantified responses using standard respiratory variables and a novel multivariate trajectory analysis. All sites recruited abdominal (ABD) expiratory activity, but rostral sites (+0.6 and +0.8~mm) produced the strongest and longest-lasting responses: ABD duration of 17.6~min versus 2.4~min at --0.2~mm ($p=0.043$, $\eta^2=0.41$), tidal volume increase of 29\% at +0.6~mm ($p=0.02$, $\eta^2=0.42$), and minute ventilation increase of 42\% ($p=0.001$, $\eta^2=0.59$). The +0.6~mm site uniquely initiated responses after the first injection and was the only location showing a post-inspiratory ABD peak. cFos analysis confirmed local activation (${\sim}$300~$\mu$m spread) without engaging PHOX2B-positive retrotrapezoid nucleus neurons above baseline. Multivariate trajectory analysis of the respiratory cycle differentiated injection sites more effectively than any single variable. These results identify the rostral pFL (+0.6~mm from VIIc) as the core site for active expiration generation.
\end{abstract}

\section{Introduction}

Breathing in mammals is generated by brainstem neural circuits comprising distinct oscillators for inspiration and expiration \citep{feldman2013breathing, smith2013brainstem}. While the inspiratory oscillator in the pre-B\"{o}tzinger complex is well-characterized, the expiratory oscillator---located in the lateral parafacial region (pFL) of the ventral medulla---remains less precisely defined \citep{janczewski2006expiratory, pagliardini2011active}. The pFL oscillator is conditionally active: it is silent at rest and engages only during metabolic demand, exercise, or pharmacological activation \citep{abdala2009abdominal, anderson2016parafacial}.

Previous studies targeting the pFL have used varied coordinates spanning approximately --0.2 to +0.5~mm from the caudal tip of the facial nucleus (VIIc), but no systematic mapping along the rostro-caudal axis has been performed \citep{huckstepp2015role, huckstepp2018chemogenetic, pisanski2020parafacial}. Chemogenetic inhibition at VIIc only partially suppressed expiratory output, and similar inhibition at +0.5~mm reduced but did not eliminate bicuculline-induced abdominal activity, suggesting that the core of the oscillator may lie at a different location.

Here we systematically map respiratory responses to bilateral bicuculline injections at five rostro-caudal positions and introduce a multivariate trajectory analysis that simultaneously assesses all respiratory signals across the full breathing cycle.

\section{Results}

\subsection{Injection Sites and Histological Verification}

Thirty-five adult male Sprague-Dawley rats ($340.4 \pm 13.2$~g) received bilateral injections of bicuculline methochloride (200~$\mu$M, 200~nL/side) or HEPES vehicle at five rostro-caudal positions relative to VIIc (Table~\ref{tab:groups}).

\begin{table}[H]
\centering
\caption{Injection group allocation.}
\label{tab:groups}
\begin{tabular}{lcccc}
\toprule
\textbf{Position (mm from VIIc)} & \textbf{n (bicuculline)} & \textbf{n (CTRL)} & \textbf{Responders} & \textbf{Post-I ABD} \\
\midrule
$-0.2$ & 5 & -- & 3/5 & No \\
$+0.1$ & 7 & -- & 7/7 & No \\
$+0.4$ & 5 & -- & 5/5 & No \\
$+0.6$ & 6 & -- & 6/6 & Yes (5/6) \\
$+0.8$ & 5 & -- & 5/5 & No \\
CTRL (HEPES) & -- & 7 & 0/7 & No \\
\bottomrule
\end{tabular}
\end{table}

cFos immunohistochemistry confirmed that bicuculline activated cells locally with an approximately 300~$\mu$m spread from the injection core. Importantly, cFos$^+$/PHOX2B$^+$ double-labeled cells did not differ between CTRL ($2.3 \pm 1.0$ per hemisection) and bicuculline groups ($2.4 \pm 0.4$; Table~\ref{tab:cfos}), confirming that the retrotrapezoid nucleus was not activated above baseline.

\begin{table}[H]
\centering
\caption{cFos cell counts at injection sites (cells per hemisection).}
\label{tab:cfos}
\begin{tabular}{lccc}
\toprule
\textbf{Group} & \textbf{Peak cFos$^+$} & \textbf{cFos$^+$/PHOX2B$^+$} & \textbf{Spread ($\mu$m)} \\
\midrule
CTRL & $44.7 \pm 4.0$ & $2.3 \pm 1.0$ & -- \\
Bicuculline (mean) & $104.5 \pm 5.1$ & $2.4 \pm 0.4$ & ${\sim}300$ \\
Beyond injection & $44.8 \pm 1.2$ & -- & -- \\
\bottomrule
\end{tabular}
\end{table}

\subsection{ABD Response Duration Varies Systematically Along the Rostro-Caudal Axis}

The most striking difference between injection sites was ABD response duration. Rostral sites (+0.6 and +0.8~mm) maintained detectable ABD activity for over 17~minutes, while the most caudal site (--0.2~mm) showed responses lasting only 2.4~minutes (Table~\ref{tab:abd_duration}).

\begin{table}[H]
\centering
\caption{ABD response temporal dynamics by injection site.}
\label{tab:abd_duration}
\begin{tabular}{lccc}
\toprule
\textbf{Position} & \textbf{Onset delay (s)} & \textbf{Duration (min)} & \textbf{Onset timing} \\
\midrule
$-0.2$ mm & $20.3 \pm 13.4$ & $2.4 \pm 1.1$ & After 2nd injection \\
$+0.1$ mm & $32.5 \pm 20.6$ & $8.2 \pm 2.4$ & After 2nd injection \\
$+0.4$ mm & $40.1 \pm 28.7$ & $10.4 \pm 3.1$ & After 2nd injection \\
$+0.6$ mm & $88.7 \pm 32.3$ & $\mathbf{17.6 \pm 2.7}$ & After 1st injection \\
$+0.8$ mm & $23.1 \pm 19.8$ & $17.1 \pm 3.3$ & After 2nd injection \\
\midrule
\multicolumn{4}{l}{ANOVA: $F(4,23) = 3.12$, $p = 0.043$, $\eta^2 = 0.41$} \\
\multicolumn{4}{l}{Tukey: $-0.2$ vs $+0.6$: $p = 0.048$; $-0.2$ vs $+0.8$: $p = 0.041$} \\
\bottomrule
\end{tabular}
\end{table}

\subsection{Tidal Volume and Minute Ventilation Increase Most at Rostral Sites}

Standard respiratory variables revealed a clear rostro-caudal gradient (Table~\ref{tab:respiratory}).

\begin{table}[H]
\centering
\caption{Peak respiratory variable changes by injection site.}
\label{tab:respiratory}
\begin{tabular}{lccccc}
\toprule
\textbf{Variable} & \textbf{$-0.2$} & \textbf{$+0.1$} & \textbf{$+0.4$} & \textbf{$+0.6$} & \textbf{$+0.8$} \\
\midrule
$\Delta$VT (\%) & $+8$ & $+16$ & $+19$ & $\mathbf{+29}$ & $+24$ \\
$\Delta$VE (\%) & $-12$ & $+6$ & $+18$ & $\mathbf{+42}$ & $+36$ \\
$\Delta$VO$_2$ (\%) & $-10$ & $-5$ & $-8$ & $-33$ & $-25$ \\
$\Delta$VE/VO$_2$ (\%) & $+2$ & $+12$ & $+28$ & $+73$ & $\mathbf{+82}$ \\
Min RR (bpm) & $37.1$ & $37.2$ & $40.1$ & $32.3$ & $30.2$ \\
\bottomrule
\end{tabular}
\end{table}

\begin{table}[H]
\centering
\caption{ANOVA results for key respiratory variables across injection sites.}
\label{tab:anova}
\begin{tabular}{lccc}
\toprule
\textbf{Variable} & \textbf{$p$-value} & \textbf{$\eta^2$} & \textbf{Best post-hoc} \\
\midrule
Max tidal volume & $0.020$ & $0.42$ & $-0.2$ vs $+0.6$: $p=0.002$ \\
Max minute ventilation & $0.001$ & $0.59$ & Rostral vs $-0.2$: $p=0.0004$ \\
Min O$_2$ consumption & $0.0001$ & $0.59$ & Rostral $<$ caudal \\
Max VE/VO$_2$ & $0.001$ & $0.55$ & Rostral $>$ caudal \\
Late-E peak amplitude & $0.41$ & -- & NS (trend toward $+0.6$) \\
\bottomrule
\end{tabular}
\end{table}

\subsection{Phase-Specific Analysis}

Kruskal-Wallis tests at individual time bins identified specific phases and time points at which injection sites differed (Table~\ref{tab:phase}).

\begin{table}[H]
\centering
\caption{Phase-specific comparisons: selected Kruskal-Wallis results.}
\label{tab:phase}
\begin{tabular}{lcccc}
\toprule
\textbf{Phase} & \textbf{Comparison} & \textbf{Time (min)} & \textbf{$H$} & \textbf{$p$ (Dunn)} \\
\midrule
Late-E airflow & $+0.6$ vs CTRL & 4 & $14.66$ & $<0.001$ \\
Late-E airflow & $+0.6$ vs CTRL & 12 & $18.35$ & $<0.001$ \\
Late-E ABD & $+0.6$ vs $-0.2$ & 8 & $11.41$ & $<0.001$ \\
Inspiratory airflow & $+0.6$ vs CTRL & 2 & $13.25$ & $<0.001$ \\
Post-I ABD & $+0.8$ vs others & 6 & $10.37$ & $0.002$ \\
\bottomrule
\end{tabular}
\end{table}

\subsection{Multivariate Trajectory Analysis}

Three-dimensional trajectories (airflow $\times$ DIA EMG $\times$ ABD EMG) computed from mean respiratory cycles revealed that rostral sites produced the largest and most sustained deviations from baseline, differentiating injection sites more effectively than any single variable (Table~\ref{tab:trajectory}).

\begin{table}[H]
\centering
\caption{Multivariate trajectory deviation from baseline by injection site.}
\label{tab:trajectory}
\begin{tabular}{lccc}
\toprule
\textbf{Position} & \textbf{Peak deviation (a.u.)} & \textbf{Significant bins} & \textbf{Duration (min)} \\
\midrule
CTRL & $0.02 \pm 0.01$ & 0 & 0 \\
$-0.2$ mm & $0.08 \pm 0.03$ & 1 & $<4$ \\
$+0.1$ mm & $0.14 \pm 0.04$ & 3 & $6$ \\
$+0.4$ mm & $0.18 \pm 0.05$ & 4 & $8$ \\
$+0.6$ mm & $\mathbf{0.38 \pm 0.06}$ & $\mathbf{8}$ & $\mathbf{16+}$ \\
$+0.8$ mm & $0.34 \pm 0.07$ & 7 & $14+$ \\
\bottomrule
\end{tabular}
\end{table}

\subsection{Inspiratory and Expiratory Duration Changes}

Across all groups, respiratory rate decreased due to prolonged expiration and shortened inspiration (Table~\ref{tab:timing}).

\begin{table}[H]
\centering
\caption{Inspiratory and expiratory duration changes (all bicuculline groups pooled).}
\label{tab:timing}
\begin{tabular}{lcccc}
\toprule
\textbf{Phase} & \textbf{Baseline (s)} & \textbf{Response (s)} & \textbf{Wilcoxon $Z$} & \textbf{$p$} \\
\midrule
Inspiratory time (Ti) & $0.318 \pm 0.043$ & $0.279 \pm 0.034$ & $3.24$ & $0.001$ \\
Expiratory time (TE) & $1.029 \pm 0.161$ & $1.313 \pm 0.188$ & $4.49$ & $0.001$ \\
\bottomrule
\end{tabular}
\end{table}

\section{Discussion}

This systematic mapping study identifies the rostral lateral parafacial region (+0.6~mm from VIIc) as the core site for active expiration generation, based on three convergent lines of evidence. First, this site produced the longest-lasting ABD responses ($>$17~min) and was the only location that initiated responses after the first injection, suggesting lowest threshold for activation. Second, it produced the largest increases in tidal volume (+29\%), minute ventilation (+42\%), and ventilatory efficiency (+73--82\%). Third, it was the only site showing a post-inspiratory ABD peak (5/6 rats), a pattern potentially reflecting engagement of post-inspiratory neurons that shape the transition between inspiration and active expiration \citep{smith2013brainstem, abdala2009abdominal}.

The multivariate trajectory analysis introduced here provides a principled approach to comparing respiratory responses across injection sites by simultaneously assessing all recorded signals across the full breathing cycle, rather than relying on individual respiratory variables that may show different site dependencies.

The confirmation that PHOX2B$^+$ RTN neurons were not activated above baseline is important, as it demonstrates that the observed respiratory effects resulted from pFL activation rather than inadvertent engagement of the adjacent chemosensitive retrotrapezoid nucleus \citep{guyenet2010retrotrapezoid, de2022phox2b}.

\section{Methods}

\subsection{Animals and Preparation}
Adult male Sprague-Dawley rats (n=35, $340 \pm 13$~g) were anesthetized with urethane (1.5--1.7~g/kg IV) after 5\% isoflurane induction. Tracheal cannulae, EMG electrodes (oblique abdominals, diaphragm), and femoral vein cannulae were placed. Bilateral vagotomy was performed. Supplemental O$_2$ (30\%) was provided.

\subsection{Injections}
Bicuculline methochloride (200~$\mu$M in HEPES with fluorobeads, 200~nL/side) was pressure-injected bilaterally at 100~nL/min using 30~$\mu$m tip pipettes. Two injections per experiment were separated by an interval. Coordinates were confirmed histologically with fluorobeads and immunohistochemistry for cFos, PHOX2B, ChAT, and TH \citep{paxinos2007rat}.

\subsection{Respiratory Measurements}
Airflow, integrated DIA and ABD EMG, and end-tidal CO$_2$ were recorded. Tidal volume, respiratory rate, minute ventilation, O$_2$ consumption \citep{depocas1957rate}, and VE/VO$_2$ were computed at 2-minute bins. Phase-specific analysis normalized signal areas for late-E, inspiratory, and post-I phases.

\subsection{Multivariate Trajectory Analysis}
Mean respiratory cycles were projected into 3D space (airflow $\times$ DIA EMG $\times$ ABD EMG). Procrustes-like distance metrics quantified deviations from baseline trajectories.

\subsection{Statistics}
One-way ANOVA with Tukey/Bonferroni post-hoc, Kruskal-Wallis with Dunn post-hoc \citep{kruskal1952nonparametric, dunn1964multiple}, Wilcoxon rank-sum tests, and repeated-measures ANOVA were used. Effect sizes: $\eta^2$. $\alpha = 0.05$.

\section{Conclusion}

Systematic rostro-caudal mapping of the lateral parafacial region identifies +0.6~mm from VIIc as the core site for active expiration generation, characterized by the strongest, longest-lasting, and most multifaceted respiratory responses to GABAergic disinhibition. Multivariate trajectory analysis provides a comprehensive framework for comparing respiratory circuit perturbations across anatomical locations.

\bibliographystyle{plainnat}
\bibliography{references}

\end{document}
